\Extrachap{Comments}
\label{chap:history}

\section*{A brief history}

Here we recall the origin of various results.

\textbf{\cref{chap:pshfunctions}}.

The notion of plurisubharmonic functions was introduced by Lelong \cite{Lel45} and Oka \cite{Oka42}, based on F.~Riesz's theory of subharmonic functions \cite{Rie26}. See \cite{Bre72} for an excellent introduction to the early history of the subject. We refer to \cite{Bre65} for the foundations of potential theory and \cite{GZ17} for the pluripotential theory.

The study of Josefson-type theorems was initiated by Josefson \cite{Jos78}, who proved \cref{thm:Josefson}.
The global Josefson theorem \cref{thm:gloJosefson} was due to Vu \cite{Vu19}. In the projective setting, it was due to Dinh--Sibony \cite{DS06} and in the K\"ahler setting, it was established by Guedj--Zeriahi \cite{GZ05}.

The extension theorem \cref{thm:GRexten} was proved in \cite{GR56}. In fact, they proved a more general version for complex spaces, see \cref{thm:GRexten2}. For some related important extension theorems, see \cite{Shi72, Wan22, CGZ13}.

\cref{prop:Kis} was due to Kiselman \cite{Kis78}.

The idea of fine topology was due to H.~Cartan \cite{Car41}.
The plurifine topology was introduced by Fuglede during the Séminaire d'analyse de Lelong--Dolbeault--Skoda of the year 1983/1984 \cite{Lel86}, but often wrongly attributed to Bedford--Taylor. The key result \cref{thm:pf_basis} was claimed in Bedford--Taylor's work \cite[Theorem~2.3]{BT87} without proof. The first rigorous proof was given by El~Marzguioui--Wiegerinck \cite{EMW06}. A weaker result was proved earlier in \cite[Theorem~4.8.7]{Kl91}.

Results in \cref{subsec:pftmfd} are certainly well-known and are already implicitly used in the literature. I could not find the proofs in the literature, so all details are presented.

The semicontinuity theorem \cref{thm:Siusemi} was due to Siu \cite{Siu74}.

Multiplier ideal sheaves were introduced by Nadel \cite{Nad90}.
The idea of \cref{thm:valuativemulti} first appeared in the ground-breaking work of Boucksom--Favre--Jonsson \cite{BFJ08}.

The openness conjecture was formulated and studied by Demailly--Koll\'ar \cite{DK01} and later proved by Berndtsson \cite{Bern13b}. The strong openness \cref{thm:strongopen} was first established by Guan--Zhou \cite{GZ15}.

\cref{lma:IandIinf} was due to \cite[Proposition~4.1.6]{Dem15}.

The notions of positivity of currents \cref{def:pos_curr} are due to Lelong \cite{Lel68} and Harvey--Knapp \cite{HK74}, although it is often mistakenly attributed to Demailly.

The $L^2$ extension theorems such as \cref{thm:OT} were first studied by Ohsawa--Takegoshi \cite{OT87}. Since then, it has become a very active area. It is almost impossible to give a complete list of the various generalizations and improvements.

\textbf{\cref{chap:npp}}

The Monge--Amp\`ere operators for bounded plurisubharmonic functions were introduced by Bedford--Taylor \cite{BT76, BT82}. The non-pluripolar product was due to Bedford--Taylor \cite{BT87}, Guedj--Zeriahi \cite{GZ07} and Boucksom--Eyssidieux--Guedj--Zeriahi \cite{BEGZ10}.

\cref{ex:1Dnpp} was claimed in \cite{BEGZ10} right after Definition~1.1 without proof. The first proof was essentially due to El Kadiri in \cite{EK24}.

The semicontinuity theorem \cref{thm:semicon} and the domination principle \cref{thm:dom_prin_1} were due to \cite{DDNL18fullmass}. The monotonicity theorem was due to \cite{WN19, DDNL18mono}.

The key lemma \cref{lma:pathoenvelope} was proved in \cite{DDNLmetric}. \cref{thm:concentration} was due to \cite{DDNLsurv}.

\textbf{\cref{chap:envelopes}}

\cref{lma:rooftopMA} was proved in \cite[Lemma~3.7]{DDNL18mono}.

The notion of the $P$-envelope is due to Ross--Witt Nystr\"om \cite{RWN14} based on the ideas of Rashkovskii--Sigurdsson \cite{RS05}.

\cref{thm:Pvarphisupport} is due to \cite[Theorem~3.8]{DDNL18mono}. The diamond inequality \cref{thm:diamond} and \cref{prop:landfinitecond1} are due to \cite[Theorem~5.4]{DDNLmetric}. Most results in \cref{subsec:fullmass} are simple generalizations of the corresponding results in \cite{DDNL18big, DDNL18fullmass}.

The study of finite energy potentials was initiated by Cegrell \cite{Ceg98} in the local setup; see also his general definition of the complex Monge--Amp\`ere operator \cite{Ceg04}. The global setup was first studied by Guedj--Zeriahi \cite{GZ07}. The variational approach to degenerate complex Monge--Amp\`ere equations in big cohomology classes was developed by Boucksom--Eyssidieux--Guedj--Zeriahi \cite{BEGZ10} and Berman--Boucksom--Guedj--Zeriahi \cite{BBGZ13}. For the complex Monge--Amp\`ere equation and pluripotential methods for weak solutions, see Ko{\l}odziej \cite{Kol98, Kol05} and the references therein.

The $\mathcal{I}$-envelope was introduced by Darvas--Xia \cite{DX22}, inspired by the works of Dano Kim \cite{Kim15} and Boucksom--Favre--Jonsson \cite{BFJ08}. The notion of $\mathcal{I}$-model singularities was first formulated in the explicit way in \cite{DX22} in 2020, although it was already essentially known in Boucksom--Jonsson's work. In fact, they correspond exactly to the homogeneous non-Archimedean potentials assuming that the relevant masses do not vanish. A less explicit equivalent formulation of $\mathcal{I}$-model potentials also appeared in \cite{Dem15}.
A few months later, the same notion was rediscovered by Trusiani \cite{Tru22}.

\cref{thm:DNTcontact} is due to \cite{DNT21}.

\textbf{\cref{chap:rays}}

The study of geodesics in the space of K\"ahler metrics was initiated by Mabuchi \cite{Mab87}, Semmes \cite{Sem92} and Donaldson \cite{Don99}. The existence of geodesics with smooth boundary data was established by Chen \cite{Chen00}.

The metric geometry and Mabuchi completion of finite-energy classes in the K\"ahler case were developed by Darvas \cite{Dar15, Dar17}. The notion of weak geodesics was studied in detail by Darvas \cite{Da17} in the K\"ahler case.

The case of general big classes was partly handled in \cite{DDNL18fullmass}, \cite{DDNL18big}.
However, the key fact that the geodesics between two full mass potentials have the correct limit at the end points does not seem to have been proved in any references. We give a proof in \cref{prop:perronenvissubgeod}. We also extend the relevant results to the relative setting.

In the special case $\phi=V_\theta$, \cref{prop:energylinear} was proved in \cite[Theorem~3.12]{DDNL18fullmass}, while the sequential case of \cref{prop:supsgeod} was proved in \cite[Proposition~3.3]{DDNL18fullmass}.

Previously, \cref{prop:geodsupsublinear} and \cref{prop:raysuplinear} were only known in the K\"ahler case.

Most results in \cref{sec:metE1} are simple extensions of \cite{DDNL18big}.

\textbf{\cref{chap:toric_ample}}

Classical toric K\"ahler geometry is rooted in the works of Atiyah \cite{Ati82}, Delzant \cite{Del88}, Guillemin \cite{Gui94}, Abreu \cite{Abr98} and Donaldson \cite{Don02}. The pluripotential-theoretic toric framework used here was first written down by Berman--Berndtsson \cite{BB13} and Coman--Guedj--Sahin--Zeriahi \cite{CGSZ19}.

The beautiful theorem \cref{thm:Yao} was first proved by Yi Yao, who did not publish the result. Later on, a new proof was found by Botero--Burgos Gil--Holmes--de Jong \cite{BBGHdJ21}. We chose to present the approach of Yao, which integrates naturally with our framework.

\Cref{prop:Etoric} is an extension of \cite[Proposition~2.9]{BB13}.

\textbf{\cref{chap:comp}}

The notion of $P$-preorder is new, as well as most results in \cref{sec:PIpartialorder}.

The $d_S$-pseudometric was introduced in \cite{DDNLmetric}. The basic properties are proved in \cite{DDNLmetric} and \cite{Xia21}.

\cref{ex:BBJ} was due to Berman--Boucksom--Jonsson \cite{BBJ21}.

\cref{thm:dSadditivity} can be viewed as a vast generalization of \cite{Cao15}, where a question raised by Demailly was confirmed.

\cref{thm:Lelongcont} is proved in \cite{Xia22}. \cref{thm:equising_cond_general} and \cref{thm:contvolu} appear to be new. These results appeared previously in the form of lecture notes.

\textbf{\cref{chap:Igood}}

The notion of $\mathcal{I}$-good singularities was due to \cite{DX21}. The name \emph{$\mathcal{I}$-good} was chosen in \cite{Xia22}.


The key results \cref{thm:charIgoodasclosure} and \cref{thm:DXmain1} are due to \cite{DX21, DX22}. Previously, Rashkovskii \cite{Ra13} also studied the local approximation problem with a different motivation.

There are some further examples of $\mathcal{I}$-good singularities provided by \cite{BBGHdJ21} with applications in the theory of modular forms in \cite{BBGHdJ22}.

\textbf{\cref{chap:trace}}

The trace operator was introduced in \cite{DX24}. Here we present a different point of view. \cref{thm:rest_volume} was proved in \cite{DX24}.

The analytic Bertini theorem \cref{thm:Bert} was proved in \cite{XiaBer}, based on the works of Matsumura--Fujino \cite{FM21} and Fujino \cite{Fuj23}. A weaker result was established by Meng--Zhou \cite{MZ23}.


\textbf{\cref{chap:testcurve}}

The technique of test curves originates from \cite{RWN14}. It was generalized by Darvas--Di Nezza--Lu \cite{DDNL18big}, \cite{DX21}, \cite{DZ22} and \cite{DXZ23}.
We give the full details of the proofs.

Test curves in \cref{def:testcur} are called \emph{maximal test curves} in the literature, a terminology which I do not like. I prefer to call the usual notion of test curves in the literature \emph{sub-test curves}.

\cref{prop:mis_RWN} was first proved by He--Testorf--Wang in \cite{HTW23}. \cref{prop:His_mass_tc} was due to Hisamoto \cite{His16}.

\cref{rmk:max_filt} was a folklore result. I am unaware of any written proof in the literature before our paper \cite{DX22}. Finski \cite{Fin22b} also gave a different proof with different techniques later on.

\cref{def:maximalgeoray} was not the original definition of maximal geodesic rays of Berman--Boucksom--Jonsson in \cite{BBJ21}. One of the first major applications of our theory was this pluripotential-theoretical characterization of maximal geodesic rays, as proved in our very first paper \cite{DX22}.

\cref{ex:RWN} is essentially due to Ross--Witt Nystr\"om \cite{RWN14}. The Phong--Sturm geodesics were first constructed in \cite{PS07}. For the related link between test configurations, filtrations and Okounkov bodies, see Boucksom--Chen \cite{BC11} and Witt Nystr\"om \cite{WN12}. For filtrations and generalized test configurations in K-stability, see Sz\'ekelyhidi \cite{Sze15}.

Results in \cref{sec:operationtc} are easy generalizations of the results proved in \cite{Xia23Operations}.

\textbf{\cref{chap:Okounkov}}

The algebraic theory of partial Okounkov bodies was developed in \cite{Xia21}.
The transcendental Okounkov body was first defined by Deng \cite{Deng17} as suggested by Demailly. The volume identity was proved in \cite{DRWNXZ}. The transcendental theory of partial Okounkov bodies is new.

\textbf{\cref{chap:bdiv}}

The notion of b-divisors was introduced by Shokurov \cite{Sho93} in the context of minimal model programs.

The applications of b-divisors in pluripotential theory began with \cite{BFJ09}. The intersection theory of nef b-divisors was introduced by Dang--Favre \cite{DF20}. The technique of singularity b-divisors was introduced in \cite{XiaPPT} in 2020. The general form first appeared in \cite{Xia22}. One year later, a special case was rediscovered in \cite{BBGHdJ21}.

The current chapter reproduces a large part of \cite{Xiabdiv}. For further developments, we refer to \cite{Xiabdiv2}.

\cref{ex:volbdivnetstrict} is reproduced from \cite[Example~3.3]{DF20}.


\textbf{\cref{chap:toric_big}}

The whole chapter appears to be new. The study of toric pluripotential theory on big line bundles was made possible by the development of partial Okounkov bodies. The key result is \cref{thm:FVtheta}.

Most results in this chapter resulted from discussions with Yi Yao.

\textbf{\cref{chap:NApotential}}

Most results in this chapter come from \cite{Xia23Operations}. Results from \cref{sec:DHmeasure} are new, although the main idea was already contained in \cite{Xia21}. The terminology of Duistermaat--Heckman measures goes back to the classical theorem of Duistermaat--Heckman \cite{DH82}; see also their addendum \cite{DH83}. In the setting of K\"ahler geometry, the Duistermaat--Heckman measures were studied in detail by Boucksom--Hisamoto--Jonsson \cite{BHJ17}.

Foundational non-Archimedean global pluripotential theory was developed by Boucksom--Favre--Jonsson \cite{BFJ16b} and, over a trivially valued field, by Boucksom--Jonsson \cite{BJ22}.

\cref{thm:DXZ} is due to \cite{DXZ23}.
An alternative approach to the transcendental theory is due to Mesquita-Piccione \cite{MP24}. For further progress in this direction, we refer to \cite{MPWN25, WN25}.

Special cases of the results in this section have been applied to study K-stability, see \cite{XiaPPT}, \cite{DZ22}, \cite{DXZ23} and \cite{DR22}. In \cite{DX22}, we established the bijective correspondence between a class of $\mathcal{I}$-model test curves and the maximal geodesic rays in the sense of \cite{BBJ21}.

\textbf{\cref{chap:Bergman}}

The study of asymptotic expansions of Bergman kernels in K\"ahler geometry was initiated by Tian \cite{Ti90} with later improvements due to Bouch\'e \cite{Bou90}, Ruan \cite{Rua98}, Catlin \cite{Cat99}, Zelditch \cite{Zel98}, Lu \cite{Lu00} among many others.

The special case of \cref{thm:pBMconvergence} without the prescribed singularity $\varphi$ was due to Berman--Boucksom--Witt Nystr\"om, see \cite{BB10}, \cite{BBWN11}. The general case is due to \cite{DX21}.
%\section*{Further results}

\section*{Open problems}
We give a list of important open problems in this theory.

We do not repeat the conjectures mentioned in the main text.

\begin{conjecture}\label{conj:exttracegeneral}
    Let $X$ be a connected compact K\"ahler manifold and $Y$ be a submanifold. Fix a K\"ahler class $\alpha$ on $X$. For each K\"ahler current $S\in \alpha|_Y$, we can find a K\"ahler current $T\in \alpha$ such that
    \[
    \Tr_Y(T)\sim_{\mathcal{I}}S.
    \]
\end{conjecture}
If we formally view $\Tr_Y$ as an analogue of the trace operator in the theory of Sobolev spaces, then this conjecture corresponds exactly to the Dirichlet problem.

\begin{problem}\label{prob:history:L209}
    Is it possible to extend the definition of the trace operator $\Tr_Y$ to the case where the ambient variety is only unibranch?
\end{problem}
The difficulty lies in the lack of Demailly type regularization theorems.


\begin{problem}\label{prob:history:L215}
    Is there a natural definition of the transcendental Okounkov body of a closed positive $(1,1)$-current $T$ with $0$-mass so that its dimension is equal to the numerical dimension of $T$?
\end{problem}
See \cite{Cao14} for the definition of the numerical dimension of a current.

The following two problems are proposed by Witt Nystr\"om.
\begin{problem}\label{prob:history:L221}
    Consider a compact K\"ahler manifold $X$ and a connected submanifold $Y$. We have defined the trace operator $\Tr_Y$ from a subset of $\QPSH(X)/\sim_\mathcal{I}$ to $\QPSH(Y)/\sim_\mathcal{I}$. Is it possible to refine this operator to one from a subset of $\QPSH(X)/\sim_P$ to $\QPSH(Y)/\sim_P$?
\end{problem}

\begin{problem}\label{prob:history:L225}
    Consider a connected compact K\"ahler manifold $X$ of dimension $n$ and a smooth flag $Y_{\bullet}$ on $X$. Consider a closed smooth real $(1,1)$-form $\theta$ on $X$ representing a big cohomology class and $\varphi\in \PSH(X,\theta)$ with $\int_X \theta_{\varphi}^n>0$.

    Can one define a refined notion of partial Okounkov bodies $\Delta'_{Y_{\bullet}}(\theta+\ddc\varphi)$ contained in $\Delta_{Y_{\bullet}}(\theta+\ddc\varphi)$ with volume given by $\frac{1}{n!}\int_X \theta_{\varphi}^n$?
\end{problem}
Note that a satisfactory solution to the latter problem is not very likely, as can be easily seen from examples on $\mathbb{P}^1$.

We also look for generalizations of our theory to more general settings.
\begin{problem}\label{prob:history:L233}
    To what extent can the results in the current book be generalized to the non-K\"ahler setting?
\end{problem}

The non-pluripolar products in the non-K\"ahler setting were recently studied by Boucksom--Guedj--Lu in \cite{BGL24}. See also the references therein.

\begin{problem}\label{prob:history:L239}
    To what extent can the results in the current book about closed positive $(1,1)$-currents be generalized to closed positive currents of higher bidegree?
\end{problem}
A fundamental issue is the lack of a strong enough Demailly type approximation for general currents. The regularization theorem of Dinh--Sibony \cite{DS04} seems too weak for our purposes.
